% Options for packages loaded elsewhere
\PassOptionsToPackage{unicode}{hyperref}
\PassOptionsToPackage{hyphens}{url}
\documentclass[
  ukrainian,
]{article}
\usepackage{xcolor}
\usepackage[margin=2.5cm]{geometry}
\usepackage{amsmath,amssymb}
\setcounter{secnumdepth}{-\maxdimen} % remove section numbering
\usepackage{iftex}
\ifPDFTeX
  \usepackage[T1]{fontenc}
  \usepackage[utf8]{inputenc}
  \usepackage{textcomp} % provide euro and other symbols
\else % if luatex or xetex
  \usepackage{unicode-math} % this also loads fontspec
  \defaultfontfeatures{Scale=MatchLowercase}
  \defaultfontfeatures[\rmfamily]{Ligatures=TeX,Scale=1}
\fi
\usepackage{lmodern}
\ifPDFTeX\else
\setmainfont{DejaVu Serif}
\setsansfont{DejaVu Sans}
\setmonofont{DejaVu Sans Mono}
\setmathfont{STIX Two Math}
\fi
\ifPDFTeX\else
  % xetex/luatex font selection
\fi
% Use upquote if available, for straight quotes in verbatim environments
\IfFileExists{upquote.sty}{\usepackage{upquote}}{}
\IfFileExists{microtype.sty}{% use microtype if available
  \usepackage[]{microtype}
  \UseMicrotypeSet[protrusion]{basicmath} % disable protrusion for tt fonts
}{}
\makeatletter
\@ifundefined{KOMAClassName}{% if non-KOMA class
  \IfFileExists{parskip.sty}{%
    \usepackage{parskip}
  }{% else
    \setlength{\parindent}{0pt}
    \setlength{\parskip}{6pt plus 2pt minus 1pt}}
}{% if KOMA class
  \KOMAoptions{parskip=half}}
\makeatother
\usepackage{graphicx}
\graphicspath{{/Users/kalinicenkovitalijmikolajovic/.codex/worktrees/5c15/Script_for_translate_book/Output_Final/}}
\setkeys{Gin}{width=0.82\linewidth,height=0.82\textheight,keepaspectratio}
\makeatletter
\newsavebox\pandoc@box
\newcommand*\pandocbounded[1]{% scales image to fit in text height/width
  \sbox\pandoc@box{#1}%
  \Gscale@div\@tempa{\textheight}{\dimexpr\ht\pandoc@box+\dp\pandoc@box\relax}%
  \Gscale@div\@tempb{\linewidth}{\wd\pandoc@box}%
  \ifdim\@tempb\p@<\@tempa\p@\let\@tempa\@tempb\fi% select the smaller of both
  \ifdim\@tempa\p@<\p@\scalebox{\@tempa}{\usebox\pandoc@box}%
  \else\usebox{\pandoc@box}%
  \fi%
}
% Set default figure placement to htbp
\def\fps@figure{htbp}
\makeatother
\ifLuaTeX
\usepackage[bidi=basic,shorthands=off,provide=*]{babel}
\else
\usepackage[bidi=default,shorthands=off,provide=*]{babel}
\fi
\ifLuaTeX
  \usepackage{selnolig} % disable illegal ligatures
\fi
\setlength{\emergencystretch}{3em} % prevent overfull lines
\providecommand{\tightlist}{%
  \setlength{\itemsep}{0pt}\setlength{\parskip}{0pt}}
\usepackage{bookmark}
\IfFileExists{xurl.sty}{\usepackage{xurl}}{} % add URL line breaks if available
\urlstyle{same}
\hypersetup{
  pdflang={uk},
  hidelinks,
  pdfcreator={LaTeX via pandoc}}

\author{}
\date{}

\begin{document}

\section{ЗАВДАННЯ}\label{ux437ux430ux432ux434ux430ux43dux43dux44f}

\begin{enumerate}
\def\labelenumi{\arabic{enumi}.}
\tightlist
\item
  Показати, що в одномірному випадку розв'язок задачі
\end{enumerate}

\[
u = u_0 + u_1 = r \cos \theta - \frac{(r^4 - 1)}{32} - \frac{(r^4 - r^2)}{24} \cos (2\theta).
\Delta u = 0, \quad 0 < r < 1 + \frac{1}{4} \sin \theta,
u \left( 1 + \frac{1}{4} \sin \theta, \ \theta \right) = \cos \theta, \quad 0 \le \theta \le 2\pi.
\begin{aligned} & (\forall \mathsf{Y}\Pi) \quad u_{tt} = c^2 u_{xx}, \\ & (\mathsf{H} \mathsf{Y}) \quad \begin{cases} u\left(x, \ 0\right) = 0, \\ u_t\left(x, \ 0\right) = \phi\left(x\right). \end{cases} \end{aligned}
u\left(1+\frac{1}{4}\sin\theta,\ \theta\right)=\cos\theta
u(1+\epsilon\sin\theta, \theta) = \cos\theta, \quad 0 \le \epsilon \le 1/4.
\]

\pandocbounded{\includegraphics[keepaspectratio]{images/_page_350_Figure_2.jpeg}}

\[
u(1+\varepsilon\sin\theta,\ \theta)=u(1,\ \theta)+u_r(1,\ \theta)(\varepsilon\sin\theta)+u_{rr}(1,\ \theta)\frac{(\varepsilon\sin\theta)^2}{21}+\ldots
\Delta u = 0, \quad 0 < r < 1 + \varepsilon \sin \theta,
u(1, \theta) + u_r(1, \theta) (\varepsilon \sin \theta) + u_{rr}(1, \theta) \frac{(\varepsilon \sin \theta)^2}{2!} + \dots = \cos \theta.
\begin{cases} u(x, 0) = x, \\ u_t(x, 0) = 0. \end{cases}
(46.8) u = u_0 + \varepsilon u_1 + \varepsilon^2 u_2 + \dots.
P_{0}
\]

\begin{enumerate}
\def\labelenumi{\arabic{enumi}.}
\setcounter{enumi}{2}
\tightlist
\item
  Проілюструвати та прокоментувати розв'язок тривимірної проблеми
\end{enumerate}

(УЧП)

\[
\begin{cases} \Delta u_{0} = 0, & 0 < r < 1 \text{ (внутри круга),} \\ u_{0}(1, \theta) = \cos \theta, & u_{0}(r, \theta) = r \cos \theta, \end{cases}
P_{1}
\]

(НУ)

\[
\begin{cases} \Delta u_{1} = 0, & 0 < r < 1 \text{ (внутри круга),} \\ u_{1}(1, \theta) = -\sin \theta \frac{\partial u_{0}(1, \theta)}{\partial r} = -\sin \theta \cos \theta \end{cases}
\]

\begin{enumerate}
\def\labelenumi{\arabic{enumi}.}
\setcounter{enumi}{4}
\tightlist
\item
  Решите аналогичную (см. задачу 3) задачу для одномерных плоских волн
\end{enumerate}

(УЧП)

\begin{itemize}
\tightlist
\item
  \begin{enumerate}
  \def\labelenumi{\arabic{enumi}.}
  \setcounter{enumi}{5}
  \tightlist
  \item
    Какова физическая природа того факта, что принцип Гюйгенса не
    справедлив в двумерном случае?
  \end{enumerate}
\item
  \begin{enumerate}
  \def\labelenumi{\arabic{enumi}.}
  \setcounter{enumi}{6}
  \tightlist
  \item
    Воспользуйтесь формулой Лейбница
  \end{enumerate}
\end{itemize}

и продифференцируйте интеграл

по времени t.

\section{СКІНЧЕННІ ПЕРЕТВОРЕННЯ ФУР'Є (ПЕРЕТВОРЕННЯ СИНУСА ТА
КОСИНУСА)}\label{ux441ux43aux456ux43dux447ux435ux43dux43dux456-ux43fux435ux440ux435ux442ux432ux43eux440ux435ux43dux43dux44f-ux444ux443ux440ux454-ux43fux435ux440ux435ux442ux432ux43eux440ux435ux43dux43dux44f-ux441ux438ux43dux443ux441ux430-ux442ux430-ux43aux43eux441ux438ux43dux443ux441ux430}

МЕТА ЛЕКЦІЇ: Ввести два нові інтегральні перетворення (скінченні
перетворення спіуса та косинуса):

\[
u_u + \frac{1}{4} u_t
\]

(конечное синус-преобразование),

\[
u_t = (1+x) u_{xx}, \\ u(x, 0) = q(x), \quad -\infty < x < \infty,
\]

(конечное косинус-преобразование),

\[
u_1 = (1 + \varepsilon_1) u_{xx}
\]

(обратное синус-преобразование),

(обратное косинус-преобразование),

та показати, як розв'язувати граничні проблеми (зокрема гетерогенні) за
допомогою цих перетворень. Раніше ми ознайомилися з регулярними
перетвореннями Фур'є та Лапласа і побачили, як застосування цих
перетворень до рівняння з частковими похідними зводить задачу до
розв'язання звичайного диференціального рівняння. Звичайне перетворення
Фур'є означає, що змінна, яку потрібно перетворити, змінюється від
\(-\infty\) до \(+\infty\) , тому воно використовується для розв'язання
задач у вільному просторі (без меж). У цій лекції ми покажемо, як
розв'язувати задачі граничних значень (з межами) шляхом перетворення
обмежених змінних (це перший раз, коли ми стикаємося з такою
трансформацією).

Давайте на деякий час забудемо про наші завдання і розберемося з самими
трансформаціями. Дамо визначення перетворень, формули циркуляції та
основні властивості. Для чого потрібні трансформації і як їх
використовувати, ми дізнаємося трохи пізніше. Однак, коротко кажучи,
суть інтегральних перетворень полягає у розкладі функцій, включених у
задачу як

частотам, решении семейства задач для каждой частоты и после-

дующем суммировании всех полученных результатов.

Почнемо з функції \(f(x)\), визначеної на інтервалі {[}0, L{]}.
Скінченні синусоїдальні та косинусоїдальні перетворення цієї функції
визначаються формулами

Конечное синус-преобразование
\(S[f] = S_n = \frac{2}{L} \int_{\sigma}^{L} f(x) \sin(n\pi x/L) dx\)
\(n = 1, 2, \dots\)
\(C[f] = C_n = \frac{2}{L} \int_{\sigma}^{L} f(x) \cos(n\pi x/L) dx\)
\(n = 0, 1, \dots\) Конечное косинус-преобразование

Читач, ймовірно, помітив, що ці перетворення не відрізняються від формул
для визначення коефіцієнтів Фур'є, коли їх розкладають на ряд синусами
або косинусами. Інверсійні формули для цих перетворень є типовими рядами
Фур'є на синусах і косинусах.

Обратное синус-преобразование

Обратное косинус-преобразование

Отметим, что в обратном косинус-преобразовании суммирование начинается с
n=0, а в синус-преобразовании --- с n=1.

\section{Приклади синусоїдального
перетворення}\label{ux43fux440ux438ux43aux43bux430ux434ux438-ux441ux438ux43dux443ux441ux43eux457ux434ux430ux43bux44cux43dux43eux433ux43e-ux43fux435ux440ux435ux442ux432ux43eux440ux435ux43dux43dux44f}

Пусть

\[
P_{0}
\]

тогда

, \((x, y, z) \in R^3\) , \(u = 0\) u\_t =

Графік функції \(f(x)\) та її перетворення показані на рис. 25.1.
Обернене перетворення записується як

\begin{enumerate}
\def\labelenumi{\arabic{enumi}.}
\setcounter{enumi}{3}
\tightlist
\item
  Решите аналогичную (см. задачу 3) задачу для двумерного уравнения
\end{enumerate}

(УЧП)

(действительная форма отображения w=z\^{}2 ).

Ця форма нам знадобиться пізніше. У цьому прикладі це цікаво, оскільки
показує, як гіперболи в площині z переходять до координатних прямих u =
const і v = const у площині w.

Тепер перейдемо до вивчення одного конкретного типу функцій комплексної
змінної, які виконують конформні відображення.

\section{Определение конформного
отображения}\label{ux43eux43fux440ux435ux434ux435ux43bux435ux43dux438ux435-ux43aux43eux43dux444ux43eux440ux43cux43dux43eux433ux43e-ux43eux442ux43eux431ux440ux430ux436ux435ux43dux438ux44f}

\begin{verbatim}
Відображення w=f(z) комплексної площини z у комплексну площину w називається конформним у точці  z_{\bullet}  площини z, якщо похідна  f'(z_{\bullet}) \neq 0  . Відображення f(z) називається конформним у регіоні D, якщо  f'(z) \neq 0  в кожній точці регіону D.
\end{verbatim}

\begin{verbatim}
Наприклад, відображення  f(z) = z^2  конформно всюди, крім точки z = 0, оскільки  z'(z) = 2z \neq 0  для всіх  z \neq 0  . З іншого боку, відображення  e^z  конформно у всій z-площині, оскільки скрізь  f'(z) = e^z \neq 0  . Яке призначення конформних відображень? Відповідь полягає в тому, що, розв'язуючи рівняння Лапласа  \phi_{xx} + \phi_{yy} = 0  в певній області площини змінних x і y, ми можемо розглядати цю площину як площину комплексної змінної z. Розглянемо комплексне відображення w = f(z) площини z у площину w, у якій введені координати (u, v). Відображуючи w = f(z  \phi_{xy} + \phi_{uy} = 0  ), рівняння Лапласа  \varphi_{xx} + \varphi_{yy} = 0  перетворюється на нове диференціальне рівняння в частинних похідних, яке залежить від нових координат і v. До — рівняння, яке  \phi_{xx} + \phi_{uu} = 0  повернеться у рівняння  \varphi_{nn} + \varphi_{nn} = 0  . Отже, ми придумаємо ідею знайти конформне відображення, яке переносить область з комплексною межею у область з простою межею (нагадаємо, що відображення  w=z^*  переносить межу першого квадранта площини z у дійсну вісь площини w).
\end{verbatim}

Ось кілька прикладів, які покажуть, як можна розв'язати диференціальні
рівняння у похідних за допомогою конформних відображень.

Уравнение Лапласа в верхней полуплоскости.

Припустимо, що нам потрібно розв'язати наступну задачу Діріхле у верхній
половині видимості (рис. 47.3):

(47.1)

Зверніть увагу, що в перших двох задачах коефіцієнти є сталими. Читач
повинен розуміти, що параметр b має бути достатньо малим, а необережний
ряд може відрізнятися

\end{document}
